% !TEX program = pdflatex
% Relatorio ABNT - DNA do Mercado: PCA e Clusterizacao
% Compilador recomendado no Overleaf: pdfLaTeX (Menu > Compiler > pdfLaTeX)
% Classe: abntex2 (disponivel na distribuicao padrao do Overleaf, nao requer instalacao extra)

\documentclass[12pt,oneside,a4paper,brazil,sumario=tradicional]{abntex2}

% ---------------------------------------------------------------------------
% PACOTES
% ---------------------------------------------------------------------------
\usepackage[T1]{fontenc}
\usepackage[utf8]{inputenc}
\usepackage[brazil]{babel}
\usepackage{times}
\usepackage{textcomp}
\usepackage{indentfirst}
\usepackage{booktabs}
\usepackage{array}

% ---------------------------------------------------------------------------
% DADOS DA CAPA E FOLHA DE ROSTO
% ---------------------------------------------------------------------------
\titulo{DNA do Mercado: PCA e Clusterização Aplicados a Ações Americanas}
\autor{Pedro Zborowski}
\local{Rio de Janeiro}
\data{2026}
\instituicao{%
  Universidade Federal do Rio de Janeiro -- UFRJ
  \par
  Computação Científica e Análise de Dados}
\tipotrabalho{Relatório técnico}
\preambulo{Relatório técnico apresentado à disciplina de Computação Científica e
Análise de Dados (UFRJ) como parte da avaliação do trabalho final, com foco na
implementação e na interpretação econômica de PCA e clusterização por K-Means.}

\begin{document}

% ---------------------------------------------------------------------------
% ELEMENTOS PRE-TEXTUAIS
% ---------------------------------------------------------------------------
\pretextual

\imprimircapa
\imprimirfolhaderosto

\begin{resumo}
Este trabalho descreve a implementação e a interpretação econômica de um pipeline
de Análise de Componentes Principais (PCA) e agrupamento por K-Means aplicado a
retornos diários de 40 ações americanas de grande capitalização, divididas a
priori em 8 setores segundo a classificação GICS. O objetivo é avaliar se a
classificação setorial de manual continua explicando o comovimento das ações ou
se o mercado as agrupa por outra lógica -- sensibilidade a juros, exposição a
commodities, exposição ao ciclo de investimento em inteligência artificial ou
risco geopolítico. Descreve-se o pipeline de dados (coleta via Yahoo Finance,
cálculo de retornos, padronização e ajuste do PCA), discute-se uma nuance técnica
da padronização utilizada -- que afeta a leitura do primeiro componente principal
como um fator de beta relativo -- e interpreta-se economicamente cada componente
principal e cada agrupamento de K-Means observado tipicamente nesse tipo de
análise. Os resultados indicam que o primeiro componente captura um fator de
mercado/beta, o segundo separa crescimento (ligado à infraestrutura de
inteligência artificial) de ativos defensivos, e o terceiro isola o setor de
energia como classe quase independente; os agrupamentos de K-Means, por sua vez,
tendem a cruzar a classificação setorial oficial em pontos economicamente
explicáveis. Conclui-se que PCA e K-Means, embora não recebam qualquer rótulo de
setor, recuperam estrutura econômica reconhecível a partir apenas dos preços,
evidenciando os limites da diversificação puramente setorial.

\vspace{\onelineskip}

\noindent
\textbf{Palavras-chave}: Análise de Componentes Principais. K-Means.
Aprendizado não supervisionado. Mercado de ações. Redução de dimensionalidade.
\end{resumo}

\pdfbookmark[0]{\contentsname}{toc}
\tableofcontents*
\cleardoublepage

% ---------------------------------------------------------------------------
% ELEMENTOS TEXTUAIS
% ---------------------------------------------------------------------------
\textual

\section{Introdução}

O aplicativo \emph{DNA do Mercado} baixa um ano de preços diários de 40 ações
americanas de grande capitalização, divididas a priori em 8 setores segundo a
classificação de manual (GICS): Tecnologia, Financeiro, Saúde, Energia, Consumo,
Industrial, Automotivo e Telecom. A partir dos retornos diários dessas ações, o
pipeline aplica duas técnicas clássicas de aprendizado de máquina não
supervisionado -- PCA (redução de dimensionalidade) e K-Means (agrupamento) --
para responder a uma pergunta específica: a classificação setorial de manual
ainda explica como essas ações se movem no dia a dia, ou o mercado agrupa as
empresas por outra lógica, como sensibilidade a juros, exposição a commodities,
exposição ao ciclo de investimento em inteligência artificial ou risco
geopolítico?

PCA e K-Means são adequados a essa pergunta porque nenhum dos dois algoritmos
recebe o rótulo de setor como entrada -- eles enxergam apenas séries numéricas de
retorno. Se os grupos encontrados pelo K-Means coincidirem com a coluna de setor,
isso é evidência de que a classificação de manual continua sendo a força
dominante. Se os grupos cruzarem setores oficiais, há uma força macroeconômica
mais forte unindo essas empresas do que a atividade declarada de cada uma.

\section{Dados e pipeline de processamento}

Os preços são obtidos via biblioteca \texttt{yfinance} (período de um ano,
preços ajustados), cobrindo cerca de 252 pregões. Ações com histórico incompleto
no período são descartadas para garantir uma matriz retangular sem lacunas. Em
seguida, os preços são convertidos em retornos percentuais diários, a unidade
correta para comparar ações de preços absolutos muito diferentes. O pipeline
segue as seguintes etapas, nesta ordem:

\begin{enumerate}
  \item coleta de preços diários ajustados na Yahoo Finance;
  \item cálculo dos retornos percentuais diários;
  \item transposição da matriz, de modo que cada empresa vire uma linha (amostra)
        e cada dia de negociação vire uma coluna (variável);
  \item padronização (\emph{StandardScaler}) coluna a coluna;
  \item ajuste do PCA sobre a matriz padronizada;
  \item agrupamento por K-Means sobre os primeiros componentes do PCA.
\end{enumerate}

Um detalhe de implementação é central para tudo o que segue: a matriz de
retornos é transposta antes da padronização e do PCA. Isso inverte o papel usual
de linhas e colunas: cada \textbf{empresa} vira uma amostra e cada \textbf{dia de
negociação} vira uma variável. Em outras palavras, o PCA não procura os dias mais
importantes do ano -- ele procura os eixos ao longo dos quais as 40 empresas mais
diferem entre si, tratando toda a trajetória de 252 dias de cada ação como sua
impressão digital.

\section{Análise de Componentes Principais (PCA)}

\subsection{Padronização e ajuste do modelo}

Antes do PCA, os dados passam pelo \emph{StandardScaler}, que padroniza cada
coluna da matriz. Como a matriz de entrada tem empresas nas linhas e dias nas
colunas, essa padronização recalcula, para cada dia individual, a média e o
desvio-padrão dos retornos das 40 empresas naquele dia, expressando o retorno de
cada empresa como um escore-$z$ em relação às demais no mesmo pregão.

\textbf{Observação técnica.} Essa padronização por dia tem uma consequência
interpretativa importante: como a média do mercado naquele dia já é subtraída
antes de o PCA rodar, o primeiro componente principal não é literalmente ``o
índice subiu ou desceu'' -- ele é o eixo que melhor separa as empresas mais
sensíveis ao mercado (beta alto) das menos sensíveis (beta baixo), porque uma
ação de beta alto se desvia da média do dia em uma direção consistente ao longo
do ano, enquanto uma ação de beta baixo quase não se desvia. Ou seja, o primeiro
componente continua sendo, na prática, um fator de beta relativo ao mercado --
apenas medido como dispersão em torno da média diária, e não como o nível do
índice em si (JOLLIFFE, 2002).

O PCA é ajustado com o número máximo de componentes permitido pelos dados (até
10, dado o tamanho do universo de ações e o número de dias disponíveis),
utilizando a implementação da biblioteca \emph{scikit-learn}
(PEDREGOSA et al., 2011). O usuário do painel escolhe, entre 2 e 10, quantos
desses componentes alimentam o K-Means; o PCA em si é recalculado apenas quando
os dados de mercado mudam.

\subsection{Interpretação econômica de cada componente}

Cada componente principal pode ser lido como um fator de risco latente -- um
padrão de comovimento que nenhuma variável do conjunto de dados original nomeia
diretamente, mas que corresponde a temas macroeconômicos reconhecíveis quando se
observa quais ações têm peso alto em cada eixo. O Quadro 1 resume a leitura
econômica típica de cada componente.

\begin{table}[htb]
\caption{Interpretação econômica dos componentes principais}
\centering
\begin{tabular}{p{1.7cm}p{2.3cm}p{9.5cm}}
\toprule
\textbf{Componente} & \textbf{Variância típica} & \textbf{Leitura econômica} \\
\midrule
PC1 & 30\%--45\% & Fator de mercado/beta. Separa ações de alta sensibilidade ao
mercado geral (tecnologia, cíclicas) das de baixa sensibilidade (defensivas). \\
PC2 & 10\%--20\% & Eixo ``crescimento/IA'' vs. ``defensivo/valor''. De um lado,
nomes ligados ao ciclo de investimento em infraestrutura de inteligência
artificial; do outro, pagadoras de dividendo estável e sensíveis a juros. \\
PC3 & 5\%--12\% & Fator commodity/energia. Isola as ações do setor de energia
como bloco quase independente do resto do mercado, reagindo a choques
geopolíticos de oferta em vez do sentimento geral de bolsa. \\
PC4--PC5 & 3\%--6\% cada & Fatores mais sutis e específicos de subgrupo (ex.:
sensibilidade de bancos a juros, exposição a tarifas de importação). \\
PC6 em diante & poucos pontos percentuais & Ruído idiossincrático -- eventos de
uma única empresa, não mais um fator de mercado amplo. \\
\bottomrule
\end{tabular}
\vspace{2mm}
\noindent{\footnotesize Fonte: elaborado pelo autor, com base nos padrões
típicos de comovimento observados em ações americanas de grande
capitalização.}
\end{table}

\subsection{Variância explicada e o número de componentes}

O painel exibe, para cada componente calculado, a fração da variância total dos
retornos que ele explica, além da soma acumulada até o número de componentes
selecionado. Em geral, a curva de variância explicada cai de forma acentuada até
o terceiro ou quarto componente e depois se achata (o chamado ``cotovelo'' do
gráfico de variância), sinalizando que os componentes seguintes agregam pouca
informação nova por unidade de complexidade.

O número de componentes usados no K-Means determina o quanto dessa estrutura é
efetivamente aproveitada, conforme resume o Quadro 2.

\begin{table}[htb]
\caption{Número de componentes principais: o que cada escolha captura}
\centering
\begin{tabular}{p{1.6cm}p{5.3cm}p{6.6cm}}
\toprule
\textbf{N\textordmasculine{} de PCs} & \textbf{O que captura} & \textbf{Consequência para os agrupamentos} \\
\midrule
2 & Apenas mercado geral (PC1) e crescimento vs. defensivo (PC2); cerca de
40\%--55\% da variância. & Separação simples de visualizar, mas energia e
industriais tendem a ficar mal classificados. \\
3--5 & Mercado, crescimento/valor, energia e um ou dois fatores sutis; cerca de
65\%--80\% da variância. & Agrupamentos interpretáveis e alinhados à intuição
econômica, cruzando o setor oficial em pontos pontuais e explicáveis. \\
8--10 & Reconstrói mais de 85\%--90\% da variância original, mas os componentes
extras já são majoritariamente ruído de empresa. & Agrupamentos mais
fragmentados; o escore de silhueta tende a cair, pois o K-Means passa a reagir a
ruído idiossincrático. \\
\bottomrule
\end{tabular}
\vspace{2mm}
\noindent{\footnotesize Fonte: elaborado pelo autor.}
\end{table}

Trata-se do trade-off clássico de viés-variância aplicado à redução de
dimensionalidade: poucos componentes dão uma visão grosseira, porém estável;
muitos componentes dão uma visão detalhada, porém ruidosa.

\section{Agrupamento por K-Means}

\subsection{Do espaço de PCA ao espaço de agrupamentos}

O K-Means não roda sobre os retornos brutos, e sim sobre as primeiras colunas
dos escores de PCA, em número igual ao escolhido pelo usuário. Essa escolha é
intencional: agrupar diretamente sobre cerca de 250 dimensões de retorno diário
seria dominado por ruído; agrupar sobre um punhado de componentes principais
concentra o algoritmo nos padrões de comovimento mais fortes já identificados
pelo PCA. O ajuste utiliza dez inicializações aleatórias diferentes, mantendo a
de menor inércia, com semente fixa para reprodutibilidade dentro da mesma sessão
de dados.

\subsection{Escolha de \emph{k} e o escore de silhueta}

Para orientar a escolha do número de agrupamentos, o painel calcula o escore de
silhueta (\emph{silhouette score}) para todo \emph{k} entre 2 e 10, usando os
mesmos componentes que alimentam o K-Means principal. O escore de silhueta mede,
em uma escala de $-1$ a $1$, o quão bem cada ponto se encaixa no seu próprio
agrupamento em relação aos agrupamentos vizinhos -- quanto maior, mais coesos e
mais separados estão os grupos.

Esse escore, no entanto, nunca deve ser o único critério de escolha. Um valor de
\emph{k} com escore ligeiramente menor, mas que produz agrupamentos
interpretáveis e alinhados a eventos reais de mercado, é mais útil para uma
análise de negócio do que o \emph{k} matematicamente ótimo que fragmenta o
resultado em grupos sem sentido econômico -- tipicamente agrupamentos de uma ou
duas ações isoladas por um evento pontual.

\subsection{Interpretação econômica por número de agrupamentos}

\begin{table}[htb]
\caption{Interpretação econômica dos agrupamentos de K-Means por valor de \emph{k}}
\centering
\begin{tabular}{p{0.9cm}p{12.3cm}}
\toprule
\textbf{\emph{k}} & \textbf{O que emerge} \\
\midrule
2 & \emph{Risk-on} vs. \emph{risk-off}: a divisão mais grosseira possível, entre
ações de alta volatilidade/beta (tecnologia, automotivas de crescimento) e ações
de baixa volatilidade (consumo básico, saúde, telecom) -- a dicotomia clássica de
mercado, redescoberta sem qualquer rótulo de setor. \\
3--4 & Ponto de equilíbrio onde o escore de silhueta costuma ser mais alto e os
agrupamentos ainda contam uma história simples: tipicamente um grupo
``IA/crescimento'', um grupo ``energia/commodities'', um grupo
``defensivas/dividendo'' e um grupo ``financeiro/cíclico''. \\
5--6 & Subestruturas dentro de setores ``artificiais'' começam a aparecer -- por
exemplo, separação entre hardware/semicondutores e software/nuvem dentro da
tecnologia, ou entre bancos de investimento e bancos comerciais tradicionais. \\
$\geq$ 7 & Superfragmentação: agrupamentos de uma ou duas ações isoladas por
eventos idiossincráticos. O escore de silhueta tende a cair -- sinal de que o
algoritmo parou de encontrar fatores de mercado amplos e passou a isolar
ruído. \\
\bottomrule
\end{tabular}
\vspace{2mm}
\noindent{\footnotesize Fonte: elaborado pelo autor.}
\end{table}

\subsection{Projeção bidimensional dos agrupamentos}

Como só é possível desenhar dois eixos por vez, o painel permite escolher quais
dois componentes visualizar em um gráfico de dispersão -- mas os agrupamentos
são sempre calculados usando todos os componentes selecionados, de modo que o
gráfico bidimensional é uma projeção parcial da separação real, que ocorre em
maior dimensão. Ao plotar, por exemplo, o primeiro componente (fator de
mercado/beta) contra o segundo (crescimento vs. defensivo), é esperado observar
os quatro agrupamentos típicos de $k=4$ ocupando regiões relativamente
compactas e distintas do plano, mesmo sabendo que parte da separação entre eles
ocorre em componentes não mostrados nesse gráfico específico.

\section{Limites metodológicos e leitura crítica}

\begin{itemize}
  \item PCA e K-Means não sabem o que é uma empresa: eles só veem números.
        Quando os agrupamentos batem com a intuição econômica (energia junta, IA
        junta), isso é evidência de que o mercado precifica essas empresas por
        fatores de risco comuns, e não um artefato do algoritmo.
  \item A classificação setorial GICS é uma simplificação: o comportamento real
        de preço frequentemente obedece a fatores macroeconômicos (juros,
        geopolítica, ciclo de investimento em IA) que atravessam múltiplos
        setores oficiais.
  \item Os resultados são dependentes do regime de mercado: como o painel
        sempre baixa um ano de dados a partir da data de execução, rodar o
        mesmo código em datas diferentes pode produzir agrupamentos diferentes
        -- isso não é uma falha, é evidência de que as correlações de mercado
        mudam conforme o contexto macroeconômico do momento.
  \item Correlação não é causalidade: o modelo mostra \emph{que} empresas se
        movem juntas; a explicação de \emph{por que} vem da leitura do
        contexto macro feita por quem interpreta os resultados, não do
        algoritmo em si.
\end{itemize}

\section{Conclusão}

O pipeline implementado demonstra, de forma reprodutível e interativa, como duas
técnicas de aprendizado não supervisionado -- PCA e K-Means -- podem extrair
estrutura econômica de dados de mercado sem qualquer rótulo de setor. O PCA
comprime cerca de 250 dimensões diárias de retorno em um punhado de fatores de
risco interpretáveis (mercado/beta, crescimento vs. valor, energia, ruído
idiossincrático); o K-Means, operando sobre esses fatores, reconstrói
agrupamentos que ora confirmam a classificação setorial de manual, ora a
contradizem de maneiras economicamente explicáveis. O resultado central deste
trabalho não é um número fixo, mas um método: uma forma sistemática de perguntar
aos próprios dados de preço quem se parece com quem, e de comparar a resposta com
o que a classificação oficial já pressupõe.

% ---------------------------------------------------------------------------
% ELEMENTOS POS-TEXTUAIS
% ---------------------------------------------------------------------------
\postextual

\phantomsection
\addcontentsline{toc}{section}{REFERÊNCIAS}
\section*{REFERÊNCIAS}

\noindent
JOLLIFFE, I. T. \emph{Principal component analysis}. 2. ed. New York: Springer,
2002.

\vspace{0.5\baselineskip}
\noindent
PEDREGOSA, F. et al. Scikit-learn: machine learning in Python. \emph{Journal of
Machine Learning Research}, v. 12, p. 2825--2830, 2011.

\vspace{0.5\baselineskip}
\noindent
YFINANCE. \emph{Documentação da biblioteca yfinance}. Disponível em:
\url{https://pypi.org/project/yfinance/}. Acesso em: 19 jul. 2026.

\end{document}
